\scanpage{155}
\mspage{71}
% archive sheet number: 154

\inserted{\deleted{\unsure{d}}c) $M$ est stable par petites $\underrightarrow{\lim}$, est accessible, \unsure{il|et} admet une filtration cardinale $(\mathrm{Filt}^{\pi}M)_{\pi\geq\pi_0}$}

\deleted{\unsure{c}}d) \deleted{\unsure{$M$}}Il existe un cardinal $\pi_{\unsure{0|1}}$, une petite catégorie $M_{\pi_0}$ stable par limites inductives de cardinal $\leq\pi_0$, \deleted{\unsure{un foncteur}} \inserted{telles qu'on} ait une équivalence de catégories

\[
M\overset{\approx}{\longleftarrow}\mathrm{Ind}_{\pi_0}\,M_{\pi_0},
\tag*{(4.3.6.1)}
\]

où le deuxième membre désigne la sous-catégorie pleine de $\mathrm{Ind}(M_{\pi_0})$ formée des \deleted{systèmes \unsure{$I$} \unsure{ind.}} ind-objets $(X_i)_{i\in I}$ de $M_{\pi_0}$ tels que $I$ soit grand devant $\pi_0$.

\uline{Démonstration de} 4.3.6. On a déjà $a'\Rightarrow a\Rightarrow b\Rightarrow a$, i.e l'équivalence \unsure{a},a$'$,b, soit $\mathrm{PTA}$ cette condition. Nous allons prouver les implications

% slanted
\margin{Pour une démonstration complète (celle donnée ici \unsure{coule|cale} à partir d'un certain moment) cf § 4.6 (\unsure{p|pp} 139--148).}

\[
\mathrm{PTA}\Longrightarrow c\Longrightarrow d\Longrightarrow\mathrm{PTA},
\]

ce qui établira le théorème.

\inserted{Prouvons} \uline{$\mathrm{PTA}\Longrightarrow c$.} Soit $S$ une petite sous-catégorie pleine \inserted{str.} génératrice. \deleted{Il est} \unsure{Ainsi|Alors}, \deleted{immédiat que si nous \unsure{a}}

